%Root main.tex
%---------------------------------------------------------------------------
%第2章　aaa
%---------------------------------------------------------------------------

\chapter{結論}
\label{chap:結論}

連系リアクトルを含まない制御モデルに対してデッドビート制御の指令値にノイズの模擬として，電流指令値の振幅に合わせた三角波を加えることで，システムの制御性がどのように変化するのかをシミュレーションによって検証した.その結果，デッドビート制御の指令値にノイズの模擬として，電流指令値の振幅に合わせた三角波を補正値として加えると特定の範囲で全高調波歪み率 (THD) の値が小さくなることを確認した.このことから，動的に指令値生成を行うことで，出力電流の制御性を改善できることが分かった.\par
今後の展望として，現在検証した範囲をより細かい範囲で検証し，制御性を向上し，制御性が向上した要因を分析する.また，事故時運転継続要件(FRT 要件)でデッドビート制御の指令値を変更することによって，システムの制御性がどのように変化するのかを検証する.さらに，繰り返し制御，AIを用いたニューラルネットワークの構築による出力電流の改善の検証を行なう．
